% titel: Blatt 10, A1: Konjugierte Paare bestimmen
% Teil II: der durchdachte Weg.

\subsection*{a) Konjugierte Basen der Säuren}
\begin{align*}
  \mathrm{HClO_4} &\to \mathrm{ClO_4^-} + \mathrm{H^+},\\
  \mathrm{H_2S}    &\to \mathrm{HS^-}  + \mathrm{H^+},\\
  \mathrm{PH_4^+}  &\to \mathrm{PH_3}  + \mathrm{H^+},\\
  \mathrm{HCO_3^-} &\to \mathrm{CO_3^{2-}} + \mathrm{H^+}.
\end{align*}

\subsection*{b) Konjugierte Säuren der Basen}
\begin{align*}
  \mathrm{Cl^-} + \mathrm{H^+} &\to \mathrm{HCl},\\
  \mathrm{SO_4^{2-}} + \mathrm{H^+} &\to \mathrm{HSO_4^-},\\
  \mathrm{H_2O} + \mathrm{H^+} &\to \mathrm{H_3O^+},\\
  \mathrm{HCO_3^-} + \mathrm{H^+} &\to \mathrm{H_2CO_3}.
\end{align*}

\begin{ergebnis}
a) ClO$_4^-$, HS$^-$, PH$_3$, CO$_3^{2-}$.
b) HCl, HSO$_4^-$, H$_3$O$^+$, H$_2$CO$_3$.
\end{ergebnis}

\begin{kontrolle}
Jedes Paar unterscheidet sich um genau ein H$^+$ und die passende Ladung
ist erfüllt (z.~B. HCO$_3^-$ $\to$ CO$_3^{2-}$). Der Ampholy-Charakter von
H$_2$O $\to$ H$_3$O$^+$ (als Base) ist korrekt.
\end{kontrolle}

\begin{merksatz}
Konjugierte Base = Minus-H$^+$, konjugierte Säure = Plus-H$^+$. Ladung
immer mitschleppen und prüfen.
\end{merksatz}
